%----------------------------------------------------------------------------------------
% Tailored CV — PhD/Postdoc, Cancer Biology & Bioinformatics
% Institute of Molecular Oncology and Functional Genomics, TUM Klinikum
%----------------------------------------------------------------------------------------

\documentclass[10pt,a4paper,sans]{moderncv}

\usepackage[utf8]{inputenc}
\usepackage[T1]{fontenc}

\moderncvstyle{classic}
\moderncvcolor{blue}

\usepackage[scale=0.78]{geometry}
\usepackage{ragged2e}

%----------------------------------------------------------------------------------------

\firstname{Kundai Farai}
\familyname{Sachikonye}

\address{Auf der Lichtung 19}{82178 Puchheim, Germany}
\mobile{(+49) 1626386344}
\email{kundai.sachikonye@bitspark.com}
\homepage{github.com/fullscreen-triangle}

%----------------------------------------------------------------------------------------

\begin{document}

\makecvtitle

%----------------------------------------------------------------------------------------
%	EDUCATION
%----------------------------------------------------------------------------------------

\section{Education}

\cventry{2019--2021}{PhD candidate, Computational Lipidomics}{Technische Universität München}{}{}{
  Doctoral research in mass spectrometry-based multi-omics, federated learning for
  distributed clinical analysis, and scientific computing pipelines.
  Departed to pursue independent computational research.
}
\cventry{2018--2019}{MSc Computational Biology}{Universität Tübingen}{}{}{}
\cventry{2014--2017}{MSc Pharmaceutical Biotechnology}{HAW Hamburg}{}{}{}
\cventry{2011--2014}{BSc Biochemistry and Cell Biology}{Jacobs University Bremen}{}{}{}

%----------------------------------------------------------------------------------------
%	RESEARCH OUTPUT
%----------------------------------------------------------------------------------------

\section{Research Output}

\cvitem{2025}{
  \textbf{On the Thermodynamic Consequences of Categorical Completion on Biological
  Systems: Specificity through VDJ Recombination Categorical Filter Cascades.}
  MHC binding affinity correlates with parent protein categorical richness
  ($r=0.73$, $p<10^{-12}$, $n=2{,}847$ IEDB peptides). AUC 0.81 vs NetMHCpan 0.68.
  Predicts collagen isoforms and tumour-associated antigens as primary immune targets
  from structural first principles. \textit{Manuscript.}
}

\cvitem{2025}{
  \textbf{On Disease Epistemology in Blind Receiver.}
  Loop-holonomy self-consistency as unique viable diagnostic criterion;
  AUC$=1.00$, 25/25 validation experiments pass. Sparse $\ell_1$ LP for
  minimum-side-effect therapeutic target selection.
  \textit{Manuscript. AIMe Registry / TUM School of Life Sciences Weihenstephan.}
}

\cvitem{2025}{
  \textbf{Syndrome: Categorical Resolution of Disease Trajectories.}
  Eight cellular oscillator classes; disease vector $\mathbf{D}\in[0,1]^8$;
  $\mathcal{O}(N\log N)$ trajectory computation. Cancer cell state transitions as
  trajectory completion in $[0,1]^8$ partition space.
  \textit{Manuscript.}
}

\cvitem{2025}{
  \textbf{Viral Infection as Host-State-Dependent Categorical Resonance.}
  Infection probability as categorical overlap integral derived from first principles.
  Oscillatory Incompleteness Theorem. \textit{Manuscript.}
}

%----------------------------------------------------------------------------------------
%	RESEARCH SOFTWARE
%----------------------------------------------------------------------------------------

\section{Research Software}

\cvitem{gospel}{
  \textbf{Large-Scale Genomic Variant Analysis and RNA-seq Integration.}
  Processes million-scale VCF files at $10^6$ variants/second; Bayesian network
  inference, Gaussian mixture models, variational Bayes for multi-omics integration.
  Validated against 1000 Genomes Phase~3, gnomAD v3.1.2, UK Biobank.
  Full pipeline: variant calling, RNA-seq quantification, pathway analysis (KEGG,
  Reactome), functional annotation, protein structure (RDKit). Deployable sequence
  homology search: \href{https://vivid-symbolism.vercel.app/search}{vivid-symbolism.vercel.app/search}.
  \href{https://github.com/fullscreen-triangle/gospel}{github.com/fullscreen-triangle/gospel}
}

\cvitem{syndrome}{
  \textbf{Cancer Cell State and Disease Trajectory Computation.}
  Eight oscillator classes (protein, enzyme, channel, membrane, ATP, genetic
  expression, calcium, circadian) with coherence indices $\eta_i\in[0,1]$;
  disease vector $\mathbf{D}\in[0,1]^8$; $\mathcal{O}(N\log N)$ trajectory
  resolution; sparse $\ell_1$ LP for minimum-intervention therapeutic target
  selection. Models cell fate transitions and metastatic progression as trajectory
  completion problems.
  \href{https://github.com/fullscreen-triangle/syndrome}{github.com/fullscreen-triangle/syndrome}
}

\cvitem{lavoisier}{
  \textbf{Multi-Omics Mass Spectrometry Pipeline.}
  Dual-pipeline: numerical MS1/MS2 analysis + visual CNN pipeline (PyTorch);
  98.9\,\% accuracy against NIST library; Ray distributed computing;
  $>$100\,GB mzML support. Virtual instrument with live forward simulation at
  \href{https://lavoisier.vercel.app/experiment}{lavoisier.vercel.app/experiment};
  Rust executable for local large-scale runs.
  \href{https://github.com/fullscreen-triangle/lavoisier}{github.com/fullscreen-triangle/lavoisier}
}

\cvitem{helicopter}{
  \textbf{Multi-Modal Microscopy Image Analysis.}
  Sequential constraint exclusion across fluorescence, brightfield, phase-contrast,
  and spectral modalities; twelve independent measurement channels; effective
  sub-nanometre structural resolution. Applicable to imaging-based CRISPR screen
  readout and tumour section analysis.
  \href{https://github.com/fullscreen-triangle/helicopter}{github.com/fullscreen-triangle/helicopter}
}

\cvitem{shakespear}{
  \textbf{Protein Analysis Instrument.}
  Folding, trajectory, catalysis, and dynamics via Kuramoto oscillator synchronization
  in S-entropy space. Applicable to oncoprotein structural analysis and cancer-associated
  protein-state transitions. Browser-based, all computation client-side. Live:
  \href{https://shakespear-nine.vercel.app/instrument}{shakespear-nine.vercel.app/instrument}.
  \href{https://github.com/fullscreen-triangle/shakespear}{github.com/fullscreen-triangle/shakespear}
}

\cvitem{honjo-masamune}{
  \textbf{Categorical Cheminformatics.}
  Zero-parameter GPU-accelerated molecular spectroscopy: vibrational modes, docking,
  3D structure, fingerprints, molecular network analysis. S-entropy coordinate framework;
  browser-side computation. Applicable to drug target characterisation and small-molecule
  analysis. Live:
  \href{https://honjo-masamune.vercel.app/tools}{honjo-masamune.vercel.app/tools}.
}

%----------------------------------------------------------------------------------------
%	RESEARCH EXPERIENCE
%----------------------------------------------------------------------------------------

\section{Research Experience}

\cventry{2019--2021}{Computational Lipidomics}{TUM / Bayerisches Forschungsinstitut}{Munich}{}{
  Mass spectrometry data pipelines (peak detection, ion annotation, feature
  engineering), multi-omics statistical modelling, federated learning for distributed
  clinical analysis. Python, Java.
}

\cventry{2017--2018}{Glycomics and Proteomics}{Universitätsklinikum Hamburg-Eppendorf}{}{}{
  Automated proteomics and glycomics characterisation pipeline for LC-MS/MS and
  MALDI-TOF data; protein identification and quantification.
}

\cventry{2019}{Bioprocess Optimisation}{Fraunhofer Institut IGB}{Stuttgart}{}{
  Real-time process monitoring and ML-based optimisation for industrial
  photobioreactors; dynamic programming, Python, C.
}

\cventry{2013}{Cancer Biomarker Discovery}{University Medical Center Groningen}{}{}{
  UPLC-MS/MS shotgun proteomics for cancer biomarker discovery; protocol
  optimisation and statistical analysis.
}

%----------------------------------------------------------------------------------------
%	TECHNICAL SKILLS
%----------------------------------------------------------------------------------------

\section{Technical Skills}

\cvitem{Genomics / Bioinformatics}{
  Variant calling (VCF, GATK), RNA-seq quantification (STAR, kallisto),
  sequence alignment (BWA, Bowtie2), single-cell analysis (Seurat/Scanpy),
  pathway analysis (KEGG, Reactome, GSEA), protein structure (RDKit, GROMACS),
  CRISPR screen analysis (MAGeCK)
}

\cvitem{ML / Statistics}{
  PyTorch, scikit-learn, XGBoost; CNNs, LSTMs, Transformers, Graph Neural Networks;
  Bayesian networks, Gaussian mixture models, variational Bayes; survival analysis;
  multi-omics dimensionality reduction (MOFA, UMAP, PCA)
}

\cvitem{Mass Spectrometry}{
  LC-MS/MS, MALDI-TOF / MALDI MSI, UPLC-MS/MS, TIMS-PASEF, QTOF;
  spectral library matching, metabolite and protein quantification
}

\cvitem{Languages}{Python, R, Rust, TypeScript, Java, C, SQL, Bash}

\cvitem{Infrastructure}{
  Distributed computing (Ray, Dask), Docker, HPC (AWS, Azure), Git,
  PostgreSQL, memory-mapped I/O for large-scale datasets, Snakemake / Nextflow
}

\cvitem{Data Formats}{
  VCF, BAM, FASTQ, FASTA, mzML, mzXML, HDF5, AnnData
}

%----------------------------------------------------------------------------------------
%	LANGUAGES
%----------------------------------------------------------------------------------------

\section{Languages}

\cvitemwithcomment{English}{Native / Fluent}{}
\cvitemwithcomment{German}{Conversational}{Living and working in Germany since 2011}

%----------------------------------------------------------------------------------------

\renewcommand{\listitemsymbol}{-~}

\end{document}
